% page 29 du fac-simile (image p-028 du scan)
% bloc 162.6 x 233.3 mm ; marge gauche 39.4 mm
\pgc{17.780mm}{0.000mm}{
\l{\cel{52}21}
\l{}
\l{\sou{alternativo}. Obligeso optar inter du selektebli. - DEFIS.}
\l{}
\l{\sou{alternatoro}.\cel{2}Elektro-mashino qua produktas elektrofluo}
\l{\cel{3}alternativa. - DEFIS.}
\l{}
\l{\sou{altitudo}. Alteso di loko, kompare a la nivelo di maro.- EFIS.}
\l{}
\l{\sou{altra}. Qua ne esas sama kam ulu, ulo, qua diferas de lu, qua}
\l{\cel{3}esas diversa. - DEFIRS.}
\l{}
\l{\sou{aludar}.\cel{2}(\sou{trans.})\cel{2}Uzar vorto, frazo, qua donas ideo pri per-}
\l{\cel{3}sono, pri kozo, pri eventajo. - EFIS.}
\l{}
\l{\sou{alumeto}.\cel{3}Stipeto de ligno, o de altra substanco, inflamebla,}
\l{\cel{3}uzata por acendar. - F.}
\l{}
\l{\sou{alumino}.\cel{2}(\sou{kemio})\cel{2}Oxido di aluminio, uzata por facar porce-}
\l{\cel{3}lano.\cel{2}Al2O3 - EFIRS.}
\l{}
\l{\sou{aluminio}. (kemio). Korpo simpla, metala, forjebla e duktila,}
\l{\cel{3}tenaca, poke densa, ofte uzata aloye kun kupro (por facar}
\l{\cel{3}juveli, vazi, e c.).\cel{2}Simbolo kemiala : Al. - DEFIRS.}
\l{}
\l{\sou{aluno}. (\sou{kemio}).So4 M2, (So4)3, M\textquotesingle{}2, 24 H2. (M : natro, o kalio,}
\l{\cel{3}amonio, rubidio, cezio, arjento ( M\textquotesingle{} = aluminio, fero,}
\l{\cel{3}magnezo, kromo o galio). - DEFIS.}
\l{}
\l{\sou{aluro}.\cel{2}La maniero marchar (pazope, trote, galope, e c.). -}
\l{\cel{3}DFR.}
\l{}
\l{\sou{aluviono}. Sedimento quan lasas la aqui kande li retroiras.-}
\l{\cel{3}DEF.}
\l{}
\l{\sou{alveolo}.\cel{2}I. Celulo vaxa quan facas la abeli. - II. Kavajo}
\l{\cel{3}di la maxil-osto en qua singla dento esas inkastrita.- DEFIS.}
\l{}
\l{\sou{amar}. (\sou{trans}.)Ligesar, per sua kordio, ad (ulu). - EFIS.}
\l{}
\l{\sou{amalgamo}. I. Aloyuro de merkurio kun metalo. - II. (\sou{metaf}.)}
\l{\cel{3}Mixuro de elementi heterogena. - DEFIRS.}
\l{}
\l{\sou{amansar}. (\sou{trans}.) I. (ulo). Igar min sovaja la animalo. -}
\l{\cel{3}II. (ulu). (\sou{metaf}.) Igar (persono) plu sociema. - IS.}
\l{}
\l{\sou{amaro}. (\sou{navig}.)Kordego uzata por retenar navo, fixigar ob-}
\l{\cel{3}jekto en navo, e c. - FIS.}
\l{}
\l{\sou{amaranto}.\cel{2}(\sou{bot}.)\cel{2}Planto autunala, kun flori reda\cel{2}bele,}
\l{\cel{3}purpuree, veluratre. - L.\cel{1}\sou{amarantus}. - DEFIRS.}
\l{}
\l{\sou{amaso}.\cel{2}To quo esas akumulita tale ke lu formacas maso,}
\l{\cel{3}granda-quanta, per adjunti intersequanta. - DEFIR.}
}
